\chapter{The laser on the cloud}
\label{ch:laser}

\devnote{This chapter describes the \code{pb-electrons} branch. The barrier and the heating are on by default with the
cloud (\opt{sfmd-laser-barrier = 1}, \opt{sfmd-laser-heating = 1}).}

The trapped electrons sit in the field of the pulse as much as the atoms do.
Two effects are modelled, both from the cycle-averaged picture of a cluster in a strong
field~\cite{ditmire1996,saalmann2006,fennel2010}: the field lowers the barrier the electrons must cross to leave, and it
heats them.
Without these the cloud is far too stable: a field-free cloud keeps $\sim$98\,\% of the electrons at every intensity,
including $10^{17}$\,W/cm$^2$, where in reality the sample is stripped (Appendix~\ref{app:validation}).

\section{The barrier}
\label{sec:barrier}

A sample of net charge $Q$ (the electrons that have left) in a uniform field $E_0$ has a potential
$-\kc Qe/r - eE_0 z$ along the field: the Coulomb well tilts and a saddle opens on the downhill side at
\begin{equation}
  r_s = \sqrt{\kc Q/E_0},\qquad \Delta W = 2\sqrt{\kc Q\,e\,E_0}
\end{equation}
below the field-free escape level (zero at infinity), with $\kc = 1.44$\,eV\,nm.
An electron less bound than $\Delta W$ can leave over the saddle.
With \opt{sfmd-laser-barrier = 1} the model uses this lowering, with the envelope amplitude $E_0(t)$ of the current
step, in three places:
\begin{description}
  \item[Escape at birth:] a new electron escapes if the depth at its atom is $D_k\le\Delta W$ (Section~\ref{sec:escape}).
  \item[The spill:] after every solve, the cloud cells whose depth $-V$ is below $\Delta W$ are emptied at once.
    Their electrons leave with their kinetic energy, $\frac32T$ each, and their electrostatic energy, computed exactly
    as $U$ before minus $U$ after; the rest of the cloud keeps its kinetic energy.
    $Q$ grows by the spilled number, so $\Delta W$ grows (as $\sqrt Q$) and the well deepens (as $Q$): the cloud is
    solved again and spills again, up to 6 rounds, until less than 1\,\% of it (or fewer than 2 electrons) would
    spill.
  \item[Evaporation:] the depth of each cell in the Maxwell-tail evaporation (Section~\ref{sec:evaporation}) is
    measured below the lowered level, $-V-\Delta W$; the leavers still take their true potential energy out of the
    cloud.
\end{description}
The cloud itself is overdense ($\omega_p$ several times $\omega$), so the laser does not reach inside it and the
equilibrium solve is unchanged; only the level the electrons must stay below moves.
The picture holds while the saddle lies outside the sample ($r_s>R$); when it does not, $\Delta W$ exceeds the well and
the cloud spills entirely, which is also what happens.
\file{laser\_log.txt} has $E_0$, $Q$, $N$, $\Delta W$, $r_s$ and the spill, escape and evaporation counts at every solve.

\paragraph{The instantaneous field.}
With \opt{sfmd-laser-barrier = 2} the spill and the evaporation level follow the instantaneous field,
$2\sqrt{\kc Q e|E(t)|}$, which breathes twice per optical cycle (2.8\,fs at 850\,nm); the solves, about 14 per cycle at
0.2\,fs, resolve it.
The escape at birth keeps $E_0$, since ADK electrons are born near the field crests.
On lysozyme mode 2 gave the same charges, ion energies (within 2\,\%) and directions as mode 1, with more spill and
less evaporation about 1\,fs earlier, but more seed noise and clamped cloud energy in 3 of 24 runs (the jumping level
lets a cold cloud's tail evaporate in bulk between solves).
Mode 1 is recommended; mode 2 is a check.

\section{Heating}
\label{sec:heating}

The cloud absorbs energy from the laser by inverse bremsstrahlung: electrons driven by the field collide with the ions
and turn the quiver motion into heat.
For a sample much smaller than the wavelength the cycle-averaged absorption is that of a small dielectric particle with
the Drude permittivity~\cite{ditmire1996},
\begin{equation}
  \varepsilon = 1 - \frac{\omega_p^2}{\omega(\omega+i\nu)},\qquad
  P = \frac{\omega}{2}\,\varepsilon_0\,\mathrm{Im}\,\varepsilon\;|E_{\mathrm{in}}|^2\,V,
\end{equation}
with the plasma frequency $\omega_p$ of the cloud's mean density $n_e = N/V$, $V = \frac43\pi R^3$, and the field inside
the cloud $E_{\mathrm{in}}$.
$P\Delta t$ is added to the cloud's kinetic and total energy every MD step (work from outside, \code{laser heating} in
the run summary and \file{laser\_log.txt}).

\paragraph{The field inside.}
For a sphere (\opt{sfmd-heating-shape = sphere}) $E_{\mathrm{in}} = 3E_0/(\varepsilon+2)$.
The formula covers the three regimes of a cluster: shielding while the cloud is dense ($\omega_p\gg\omega$, small
$E_{\mathrm{in}}$), the Mie resonance as it expands through $\omega_p = \sqrt3\,\omega$ (a density of $3n_c$,
$n_c\approx1.55\times10^{21}$\,cm$^{-3}$ at 850\,nm), and the resistive regime $\nu\gg\omega$.
With \opt{sfmd-heating-shape = ellipsoid} (the default) the cloud has the shape of the sample's gyration ellipsoid,
semi-axes $a_i$, and the field along each axis is $E_{0,i}/(1+L_i(\varepsilon-1))$, with the depolarization factors
\begin{equation}
  L_i = \frac{a_1a_2a_3}{2}\int_0^\infty\frac{\dd s}{(s+a_i^2)\sqrt{(s+a_1^2)(s+a_2^2)(s+a_3^2)}}
      = \frac{a_1a_2a_3}{3}\,R_D(a_j^2,a_k^2,a_i^2)
\end{equation}
(Carlson's elliptic integral; $\sum L_i = 1$, $L_i = \frac13$ for a sphere).
The model has no field direction, so it averages over the orientation of the polarization:
\begin{equation}
  |E_{\mathrm{in}}|^2 = \frac{E_0^2}{3}\sum_i\big|1+L_i(\varepsilon-1)\big|^{-2}.
\end{equation}
The resonance then falls at $n_c/L_i$ for each axis rather than at $3n_c$; lysozyme starts at $L$ = 0.21, 0.37, 0.41.
On lysozyme the ellipsoid gave 13\,\% more heating and 6\,\% more ion energy at $10^{15}$\,W/cm$^2$ and about 1\,\% at
$10^{16}$, directions within the noise.

\paragraph{The collision rate.}
$\nu$ is the electron--ion collision rate~\cite{nrl},
\begin{equation}
  \nu_{ei} = 2.91\times10^{-6}\,Z\,n_e[\mathrm{cm^{-3}}]\,\ln\Lambda\,T^{-3/2}\ \mathrm{s^{-1}},\qquad
  Z = \frac{\sum q_a^2}{\sum q_a},\qquad \ln\Lambda = \max\big(2,\ 24-\ln(\sqrt{n_e}/T)\big),
\end{equation}
at the cloud's temperature (at least the floor, Section~\ref{sec:temperature}), with two corrections:
\begin{description}
  \item[The cap] $\nu_{ei}\le$ \opt{sfmd-heating-nu-cap}$\cdot\omega_p$ (1): a cold, dense cloud is strongly coupled, and
    the weak-coupling formula would give $\nu\gg\omega_p$; a collision rate saturates near the plasma frequency.
    The cap is an order-of-magnitude choice, not fitted.
    It matters only while the cloud is cold and dense: on lysozyme at $10^{15}$\,W/cm$^2$, a cap of 0.5 gave 16\,\% less
    heating and 8\,\% less ion energy; at $10^{16}$ nothing.
    No cap gave the same as cap 1.
    0 removes it.
  \item[Surface damping] $\nu \mathrel{+}= A\,\bar v/R$ with $\bar v = \sqrt{8T/\pi m_e}$ and $A$ =
    \opt{sfmd-heating-surface} (1): the electrons of a nm-sized cloud cross it in $R/\bar v$, which damps its plasmon
    (Kreibig damping~\cite{kreibig1995}).
    Without it only $\nu_{ei}$ damps the resonance, and a hot, thin cloud ($\nu/\omega\sim10^{-5}$) can absorb MeV
    within a step when a solve lands on the resonance: 15--19\,MeV heating outliers in about 1 of 8 runs at
    $10^{16}$--$10^{17}$\,W/cm$^2$.
    With $A=1$ ($\nu_s/\omega\approx$ 0.5--0.9 near the resonance) the outliers are gone and the seed spread of the ion
    energy falls from $\pm90$ to $\pm28$\,keV at $10^{16}$.
    $A$ is of order 1 in the literature and not fitted; from $A$ = 0.5 to 2 the ion energy at $10^{16}$ rises by 2.3\,\%,
    directions unchanged.
\end{description}

\section{Linear polarization}
\label{sec:polarization}

\devnote{\opt{sfmd-laser-polarization} is new (tpr 185) and not validated beyond the checks below. Default 0: the laser
without a direction, as in the rest of this chapter.}

The barrier above is lowered all round, which overstates escape: the real saddle exists only on the downhill side of
the field, on the polarization axis, and moves from one pole to the other every half cycle.
With \opt{sfmd-laser-polarization = x}, \code{y} or \code{z} the field $E(t)\,\hat e$ of the pulse (carrier included)
acts on the cloud and the ions directly:
\begin{description}
  \item[In the solve.] The laser's potential energy of an electron, $U_L = eE(t)\,(\mathbf r-\mathbf c)\cdot\hat e$
    ($\mathbf c$ the heavy-atom centre), is added to the potential of the PB solve, $V = -\Phi + Kn + U_L$, so the cloud
    polarizes: it shifts along $\hat s = -\mathrm{sign}(E)\,\hat e$, the direction the field pushes electrons.
    The cloud's plasma frequency is several times the laser's, so it follows the field adiabatically; the solves sample
    the carrier (14 per cycle at 0.2\,fs and 850\,nm; \code{grompp} warns below 8) and the field is held in between.
  \item[The escape level] is the saddle on the downhill side. Along $\hat s$, outside the sample,
    \begin{equation}
      U(r) = -\frac{\kc Q}{r} - \frac{\kc P}{r^2} - |E|\,r,\qquad P = \mathbf p\cdot\hat s,
    \end{equation}
    with $\mathbf p$ the dipole of ions and cloud about $\mathbf c$ from the last solve (the shifted cloud gives $P<0$,
    which raises the saddle). The saddle $r_s$ is the largest root of $|E|r^3-\kc Qr-2\kc P=0$, the level
    $U_s = U(r_s)$; for $P=0$ it is the barrier of Section~\ref{sec:barrier}, $U_s=-\Delta W$.
    If that root lies inside the sample ($r_s<R_e$, the gyration ellipsoid's radius along the axis), or there is
    none because the dipole term wins (a strongly polarized cloud of small $Q$), $U$ falls all the way out beyond
    $R_e$ and the level is $U(R_e)$, the highest point on the way out. (The first version took the root, or for no
    root a point deep inside the sample, where the dipole term made the level huge; a neutral cloud along the long
    axis was then held at $10^{17}$\,W/cm$^2$ and absorbed 130\,MeV in one of 8 runs.) Beyond the saddle plane
    the field term is held at its value there: those cells are outside the trap.
  \item[Escape] is then decided by energy against $U_s$: a cell spills if $V\ge U_s$, a newborn electron escapes if
    $-D_k + U_L(\mathbf x_k)\ge U_s$, and evaporation measures each cell's depth as $U_s - V$. Electrons on the
    up-field side are lifted by $U_L$, those on the down-field side lowered; the cloud's tail beyond the saddle region
    goes first. Both \opt{sfmd-laser-barrier} modes behave the same with the polarization (it needs one of them on).
  \item[The ions] feel $qE\,\hat e$ of the solve's field, so that ions and cloud together feel the net charge's
    $QeE$.
  \item[Heating] with the ellipsoid uses the field along the axis, $|E_{\mathrm{in}}|^2 = E_0^2\sum_i (\hat e\cdot\hat u_i)^2
    |1+L_i(\varepsilon-1)|^{-2}$ with the sample's principal axes $\hat u_i$, instead of the orientation average.
\end{description}

\paragraph{Energy.}
The cloud's energy ledger is kept without the field term ($E_{\mathrm{kin}} = E_{\mathrm{tot}} - U$, $U$ the
electrostatic energy of cloud and ions), and the field's work on the cloud between two solves,
$\int E\,\dd X$ with $X = \sum_i n_i\,\mathbf r_i\cdot\hat e$, is added to $E_{\mathrm{tot}}$ with the mean of the two
solves' fields. For a cloud that follows the field linearly this sum telescopes to zero over a cycle; with the new
field alone it would put $\frac12\alpha\,\Delta E^2$ into the cloud at every solve, several keV per solve for lysozyme at
$10^{16}$\,W/cm$^2$. Births, captures, evaporation and spills update $X$ with the electrons' positions, so the work
done on an electron born at its atom as it joins the shifted cloud is counted. The run's end reports the cumulative
polarization work. The part stored in the shifted cloud, $\frac12\alpha E^2$, returns as the field falls; electrons born at
their atoms into the shifted cloud gain about $|E|d$ each ($d$ the shift), and the down-field tail that spills was
pulled there by the field, so a cloud that changes during the pulse keeps a net gain (643\,keV for lysozyme at
$10^{16}$\,W/cm$^2$, against 1.6\,MeV of heating); \file{laser\_log.txt} gets the solve's field,
$p_{\hat e}$, the work and the solves without a saddle as extra columns ($\Delta W$ and $r_s$ there are $-U_s$ and
the saddle's distance).

\section{What is left out}

\begin{itemize}
  \item The field's direction, unless \opt{sfmd-laser-polarization} is set: the barrier is lowered all round, not
    only along the polarization, and the heating is orientation-averaged.
  \item Quiver motion and the electrons' response within a cycle (except the spill with mode 2).
  \item Heating at the surface (vacuum or Brunel heating), and the density profile: one mean density.
  \item The energy of electrons that leave: they take their thermal energy, not a drift energy from the field.
\end{itemize}
